% ProtoCal -- manuscript boilerplate
%
% Target venues:
%   1. arXiv (cs.LG / cond-mat.mtrl-sci) -- initial preprint
%   2. AI for Science, IOP Publishing (ISSN 3050-287X), gold open access.
%      Interdisciplinary, AI applied to scientific discovery. Submission via
%      ScholarOne. CRediT contribution statement encouraged.
%      https://iopscience.iop.org/journal/3050-287X/page/author-guidelines
%   3. Materials Letters (Elsevier) -- short-format fallback. Strict limit
%      (~2500 words / 4 printed pages incl. figures). If targeting this,
%      cut to Findings 1 and 2 only and drop the synthetic-recovery section.
%
% Build: pdflatex main && bibtex main && pdflatex main && pdflatex main

\documentclass[11pt]{article}
\usepackage[utf8]{inputenc}
\usepackage[T1]{fontenc}
\usepackage{amsmath,amssymb}
\usepackage{booktabs}
\usepackage{graphicx}
\usepackage{hyperref}
\usepackage[margin=1in]{geometry}
\usepackage{natbib}

\title{Protocol-Aware Hierarchical Modelling of Scattered Literature Data\\
       for Dental Resin Composites}
\author{Shrey Patel\thanks{Corresponding author.}}
\date{\today}

\begin{document}
\maketitle

\begin{abstract}
Machine learning for dental biomaterials is limited less by data volume than by
uncontrolled heterogeneity: published studies differ in specimen geometry, storage
medium, aging regime and curing protocol, and these nuisance factors are entangled with
composition. We formulate a hierarchical model that conditions explicitly on reported
protocol, gauge-fixed at a reference protocol so that the protocol term and per-study
offsets are identifiable when each study reports under a single protocol. Using 1{,}802
specimen-level measurements spanning 13 materials, 4 aging protocols and 5 properties,
we show that protocol accounts for up to 32\% of variance in flexural strength and
Vickers hardness. We further show, by variance decomposition, that an \emph{additive}
protocol term is mis-specified for flexural strength, where the material\,$\times$\,protocol
interaction exceeds measurement noise by a factor of 1.42. Finally, synthetic-recovery
experiments establish that the shared protocol function is recoverable at realistic
literature sparsity ($r \approx 0.95$) while per-study offsets are not ($r \approx 0.31$),
bounding what can honestly be claimed from literature-derived corpora.
\end{abstract}

\section{Introduction}
% Motivate: small, heterogeneous, inconsistently reported datasets are THE bottleneck.
% Existing dental-composite ML pools multi-study tables and validates with random k-fold.
% Cite Paniagua 2025; note no public literature-extracted dataset with protocol metadata.
\textbf{[TODO]}

\section{Related work}
% - ML for dental composites: \citep{paniagua2025}
% - active learning / BO for dental materials
% - hierarchical / mixed-effects meta-analysis; domain generalisation (DANN)
% - leave-one-cluster-out evaluation in materials informatics
\textbf{[TODO]}

\section{Method}
For a property value $y$ of formulation $i$ reported by study $s$,
\begin{equation}
  y_{i,s} = f_{\text{shared}}(x_i) + g_{\text{proto}}(p_s) + b_s + \varepsilon_s ,
  \label{eq:model}
\end{equation}
where $x_i$ is the composition/process vector, $p_s$ the protocol vector, $b_s$ a
per-study offset under a learned $\mathcal{N}(0,\tau^2)$ prior and $\varepsilon_s$
per-study heteroscedastic noise.

\paragraph{Identifiability.}
Because a study typically reports under one fixed protocol, $g_{\text{proto}}(p_s)$ and
$b_s$ are both per-study additive constants and are not separable in isolation. We
resolve this structurally rather than by assumption:
\begin{enumerate}
  \item \textbf{Gauge fixing.} We define
        $\tilde g(p) = g(p) - g(p_{\text{ref}})$, so $\tilde g(p_{\text{ref}}) \equiv 0$
        by construction, removing the additive ambiguity and making
        $f_{\text{shared}}$ interpretable as the reference-protocol prediction.
  \item \textbf{Asymmetric regularisation.} $g_{\text{proto}}$ is low-capacity, shared
        across studies and heavily weight-decayed; $b_s$ is a free scalar per study
        under strong shrinkage. Between-study variation explainable by observed
        protocol covariates is therefore cheaper to place in $g_{\text{proto}}$.
  \item \textbf{Synthetic recovery as a precondition.} Section~\ref{sec:recovery}.
  \item \textbf{Leakage diagnostic.} Regressing fitted $b_s$ on $p_s$ post hoc;
        $R^2 \approx 0$ confirms no protocol signal leaked into the offsets.
\end{enumerate}

\section{Data}
\textbf{[TODO]} Lodz aging corpus (Zenodo 10.5281/zenodo.14576242, CC-BY-4.0):
1{,}802 specimen measurements, 13 materials, 4 protocols, 5 properties, 238/260 cells
(91.5\% complete factorial). Protocol encoded per file and parsed into medium, duration
and temperature.

\section{Results}

\subsection{Protocol is a first-order effect}
\begin{table}[h]\centering
\begin{tabular}{lrrrr}
\toprule
Property & Material & Protocol & Interaction & Inter./noise \\
\midrule
Flexural strength   & 55.5\% & 31.2\% & 13.3\% & 1.42 \\
Vickers hardness    & 49.1\% & 32.3\% & 18.5\% & 0.78 \\
Diametral tensile   & 60.4\% & 19.8\% & 19.8\% & 0.51 \\
Fracture toughness  & 79.5\% & 11.6\% &  8.9\% & 0.30 \\
Flexural modulus    & 93.6\% &  2.5\% &  3.9\% & 0.23 \\
\bottomrule
\end{tabular}
\caption{Variance decomposition on log values, material $\times$ protocol.}
\label{tab:variance}
\end{table}

\subsection{The additive protocol term is mis-specified for strength}
\textbf{[TODO]} Interaction exceeds within-group noise for flexural strength
(1.42$\times$); propose low-rank interaction $g(p) + \langle u(x), v(p)\rangle$.

\subsection{What is recoverable}
\label{sec:recovery}
\begin{table}[h]\centering
\begin{tabular}{lrrr}
\toprule
Setting & Obs/(study,prop.) & $g_{\text{proto}}$ & $b_s$ \\
\midrule
Sparse (as published) &  6.2 & 0.951 & 0.313 \\
$2\times$ rows        & 14.3 & 0.950 & 0.435 \\
Full reporting        & 13.5 & 0.957 & 0.525 \\
$4\times$ + full      & 59.4 & 0.968 & 0.558 \\
\bottomrule
\end{tabular}
\caption{Synthetic recovery (Pearson $r$ against ground truth).}
\label{tab:recovery}
\end{table}

\section{Limitations}
Per-study offsets are not reliably estimable at literature density, so $b_s$ must not be
read as a study-quality signal. The corpus analysed here is single-laboratory; the
protocol axis is aging regime only, not curing or geometry. No patient-level data is
used or required at any stage.

\section{Conclusion}
\textbf{[TODO]}

\bibliographystyle{plainnat}
\bibliography{refs}
\end{document}
